\documentclass[UTF8]{ctexart}
\usepackage{geometry,booktabs,longtable,array,tabularx}
\geometry{a4paper,margin=2.0cm}
\setlength{\emergencystretch}{3em}
\setlength{\tabcolsep}{3pt}
\renewcommand{\arraystretch}{1.12}
\title{JiuwenSwarm v4 可审计科研论文半流程报告}
\author{JiuwenSwarm/UCR}
\date{}
\begin{document}
\maketitle
\begin{abstract}
本文汇总三个主题的科研半流程演示。系统检索公开资料元数据，执行 JiuwenSwarm/UCR 本地证据实验，生成研究计划和报告，并保存 Claim--Evidence 账本。本文不把流程基准包装成领域科学结论；正式引用、实验真实性、研究结论和最终提交资格仍需人工确认。
\end{abstract}
\section{统一边界}
UCR 的标签是 \texttt{process\_evidence\_only}，它表示系统的流程证据和审计能力，不等于上下文工程、记忆引擎或自我进化的领域效果。所有来源默认处于待审核状态。当前 PDF 是由本机 XeLaTeX 编译生成的正式查看产物，源文件和实验记录同时保留在提交包中。
\section{矩阵概览}
\begin{longtable}{@{}>{\raggedright\arraybackslash}p{0.30\textwidth}>{\raggedright\arraybackslash}p{0.42\textwidth}>{\raggedleft\arraybackslash}p{0.14\textwidth}@{}}
\toprule
主题 & 状态 & Provider 调用\\\\
\midrule
\endfirsthead
\toprule
主题 & 状态 & Provider 调用\\\\
\midrule
\endhead
context-\allowbreak{}engineering & \texttt{\scriptsize completed\_\allowbreak{}pending\_\allowbreak{}human\_\allowbreak{}review} & 3 \\
memory-\allowbreak{}engine & \texttt{\scriptsize completed\_\allowbreak{}pending\_\allowbreak{}human\_\allowbreak{}review} & 3 \\
self-\allowbreak{}evolution & \texttt{\scriptsize completed\_\allowbreak{}pending\_\allowbreak{}human\_\allowbreak{}review} & 3 \\
\bottomrule
\end{longtable}
\section{Agent 上下文工程}
\textbf{摘要：}本报告检验程序级证据护栏能否提高 Agent 执行报告的证据一致性。三个 arm（no-\allowbreak{}rail、prompt-\allowbreak{}only、full-\allowbreak{}rail）均完成且 return\_\allowbreak{}code=0，task\_\allowbreak{}success\_\allowbreak{}rate 均为 1.0。no-\allowbreak{}rail 的 UCR=0.5556（5 / 9），prompt-\allowbreak{}only 与 full-\allowbreak{}rail 的 UCR 均为 0.0（0 / 4）。full-\allowbreak{}rail 的 UCR 低于 no-\allowbreak{}rail，但与 prompt-\allowbreak{}only 相同，存在天花板效应可能。UCR 仅为过程证据，不推断任务正确性或因果效应。\\
\textbf{结论：}在本次实验范围内，full-\allowbreak{}rail 的 UCR（0.0）低于 no-\allowbreak{}rail（0.5556），与 prompt-\allowbreak{}only（0.0）相同，部分支持假设但无法区分 full-\allowbreak{}rail 与 prompt-\allowbreak{}only 的效应。所有结论需绑定 claim\_\allowbreak{}id 与实验记录或来源引用，且需人工确认 UCR 的解释边界。\\
\textbf{限制：}UCR 仅为过程证据，不推断任务正确性或因果效应。CFR 与 NFR 未测量，不能用于支持或反驳假设。full-\allowbreak{}rail 与 prompt-\allowbreak{}only 的 UCR 均为 0.0，需人工判断是否存在天花板效应。citation\_\allowbreak{}allowed=false 的材料不得作为正式引用，citation\_\allowbreak{}coverage=0.0。claim-\allowbreak{}002、claim-\allowbreak{}003、claim-\allowbreak{}005、claim-\allowbreak{}007、claim-\allowbreak{}009、claim-\allowbreak{}018、claim-\allowbreak{}019、claim-\allowbreak{}020 为 unsupported\_\allowbreak{}claims，需人工复核。
\section{Agent 记忆引擎}
\textbf{摘要：}本报告检验结构化记忆记录是否有助于 Agent 保持审计上下文一致。三个 arm（no-\allowbreak{}rail、prompt-\allowbreak{}only、full-\allowbreak{}rail）均完成且 return\_\allowbreak{}code=0，任务成功率均为 1.0。UCR 在 no-\allowbreak{}rail 为 0.5556（5 / 9），在 prompt-\allowbreak{}only 与 full-\allowbreak{}rail 均为 0.0（0 / 4）。UCR 仅作为过程证据解读，不外推领域效果。\\
\textbf{结论：}结构化记忆记录可能改善流程可追溯性，但当前基准不证明领域效果。no-\allowbreak{}rail 的 UCR 为 0.5556（claim-\allowbreak{}001、claim-\allowbreak{}004、claim-\allowbreak{}006、claim-\allowbreak{}008），prompt-\allowbreak{}only 与 full-\allowbreak{}rail 的 UCR 均为 0.0（claim-\allowbreak{}010 至 claim-\allowbreak{}017）。结论需人工确认，且不得将 UCR 外推为领域效果。\\
\textbf{限制：}UCR 仅为 process\_\allowbreak{}evidence\_\allowbreak{}only，不证明领域效果。CFR 与 NFR 未测量。所有材料 citation\_\allowbreak{}allowed=false，不得作为正式引用（claim-\allowbreak{}018、claim-\allowbreak{}019、claim-\allowbreak{}020）。citation\_\allowbreak{}coverage=0.0，allow\_\allowbreak{}paper=false。claim-\allowbreak{}002、claim-\allowbreak{}003、claim-\allowbreak{}005、claim-\allowbreak{}007、claim-\allowbreak{}009、claim-\allowbreak{}018、claim-\allowbreak{}019、claim-\allowbreak{}020 需人工复核。
\section{Agent 自我进化}
\textbf{摘要：}本报告检验带证据的反思循环能否避免把未执行任务写成已完成。三个实验臂均完成且返回码为0：no-\allowbreak{}rail 的 UCR=0.5556（5 / 9），prompt-\allowbreak{}only 与 full-\allowbreak{}rail 的 UCR=0.0（0 / 4）。UCR 仅作为过程证据，不证明真实自我进化。\\
\textbf{结论：}当前基准下，带证据的反思循环与审计表达改善相关：no-\allowbreak{}rail 出现 5 个 unexecuted，而 prompt-\allowbreak{}only 与 full-\allowbreak{}rail 均为 0。但该结果不证明真实自我进化，也不建立因果作用；需补充独立基准并完成人工确认。\\
\textbf{限制：}UCR 仅为过程证据，不能推断真实自我进化或泛化能力。CFR 与 NFR 未测量。citation\_\allowbreak{}allowed=false 的材料不得作为正式引用，citation\_\allowbreak{}coverage=0.0。claim-\allowbreak{}002、claim-\allowbreak{}003、claim-\allowbreak{}005、claim-\allowbreak{}007、claim-\allowbreak{}009、claim-\allowbreak{}018、claim-\allowbreak{}019、claim-\allowbreak{}020 仍待人工复核。full-\allowbreak{}rail 的 rail.attempts=1 与 accepted=true 是否足以支持护栏有效性仍需确认。
\section{编译与审核说明}
该统一报告和三个主题报告都需要人工阅读。编译命令为：
\begin{verbatim}
xelatex -interaction=nonstopmode -halt-on-error paper.tex
xelatex -interaction=nonstopmode -halt-on-error paper.tex
\end{verbatim}
最终状态保持为 \texttt{prepared\_not\_uploaded}。
\end{document}
